% GENERATED FILE. Edit canonical lessons, then run npm run generate:books.
% canonical-source-hash: fnv1a64:c115921b5fb3f6c5
% canonical-lessons: BN-C418-munapha, BN-C418-shulko, BN-C418-khajna, BN-C418-mohajon, BN-C418-baki

\chapter{Profit, Taxes and Moneylenders}
\label{ch:bn-profit-taxes-and-moneylenders}
\hlchaptermodality{418}

\begin{chapteropening}
\textbf{By the end of this chapter:} \emph{I can talk about profit, duty and land tax, a moneylender and what is left to pay.}

Complete the last lesson of chapter 418: I can talk about profit, duty and land tax, a moneylender and what is left to pay.
\end{chapteropening}

\section[\texorpdfstring{\bn{মুনাফা} (munāphā)}{মুনাফা (munāphā)}]{\bn{মুনাফা} (munāphā) — profit (in trade)}

\label{lesson:BN-C418-munapha}



\begin{quote}
Before the new one: say the Bengali for to transfer (money), then the Bengali for a contribution (to a fund).
\end{quote}

\subsection*{You'll want to know: \bn{মুনাফা}}
\textbf{\bn{মুনাফা}} — \emph{munāphā} — \textquotedblleft{}profit (in trade)\textquotedblright{}.

\begin{quote}
\bn{এই} \bn{বছর} \bn{ব্যবসায়} \bn{ভালো} \bn{মুনাফা} \bn{হল।}
\end{quote}

\begin{quote}
\emph{ei bôchhôr bæbsāy bhālo munāphā hôlô.}
\end{quote}

\begin{quote}
The business made a good profit this year.
\end{quote}

\bn{লাভ}, gain, is the wider word; \bn{মুনাফা} is the trader's profit.

\begin{cousinweb}
\textbf{\bn{মুনাফা}} (\emph{munāphā}) is usually traced to Arabic \emph{manfaʿa}, \textquotedblleft{}use, benefit, profit\textquotedblright{}, which reached Bengali through Persian.
\end{cousinweb}

\begin{grammarlens}[title={What you've built}]
One more word for getting around.
\end{grammarlens}

\subsection*{Your turn}
\begin{itemize}
\raggedright
  \item \emph{Say it:} \emph{munāphā}
  \item \emph{Say it:} \emph{munāphā}, once more
  \item \emph{Say it:} \emph{ei bôchhôr bæbsāy bhālo munāphā hôlô.} — The business made a good profit this year.
  \item \emph{Recall:} say the Bengali for to transfer (money), then the Bengali for a contribution (to a fund), then say \emph{munāphā} again
\end{itemize}

\begin{tcolorbox}[breakable,colback=teal!4,colframe=teal!35!black,title={Before you move on}]
What does \bn{মুনাফা} mean? (\textquotedblleft{}Profit (in trade)\textquotedblright{}.) Say it once more.
\end{tcolorbox}
\section[\texorpdfstring{\bn{শুল্ক} (shulkô)}{শুল্ক (shulkô)}]{\bn{শুল্ক} (shulkô) — a duty; a tax on goods}

\label{lesson:BN-C418-shulko}



\begin{quote}
Before the new one: say the Bengali for a contribution (to a fund), then the Bengali for profit (in trade).
\end{quote}

\subsection*{You'll want to know: \bn{শুল্ক}}
\textbf{\bn{শুল্ক}} — \emph{shulkô} — \textquotedblleft{}a duty; a tax on goods\textquotedblright{}.

\begin{quote}
\bn{বিদেশি} \bn{জিনিসের} \bn{ওপর} \bn{শুল্ক} \bn{দিতে} \bn{হয়।}
\end{quote}

\begin{quote}
\emph{bideshi jiniser opôr shulkô dite hôy.}
\end{quote}

\begin{quote}
You have to pay duty on foreign goods.
\end{quote}

\begin{cousinweb}
\textbf{\bn{শুল্ক}} (\emph{shulkô}) is Sanskrit \emph{śulka}, \textquotedblleft{}a toll, a tax, a duty, especially at ferries and roads\textquotedblright{} (Monier-Williams), taken over whole.
\end{cousinweb}

\begin{grammarlens}[title={What you've built}]
One more word for getting around.
\end{grammarlens}

\subsection*{Your turn}
\begin{itemize}
\raggedright
  \item \emph{Say it:} \emph{shulkô}
  \item \emph{Say it:} \emph{shulkô}, once more
  \item \emph{Say it:} \emph{bideshi jiniser opôr shulkô dite hôy.} — You have to pay duty on foreign goods.
  \item \emph{Recall:} say the Bengali for a contribution (to a fund), then the Bengali for profit (in trade), then say \emph{shulkô} again
\end{itemize}

\begin{tcolorbox}[breakable,colback=teal!4,colframe=teal!35!black,title={Before you move on}]
What does \bn{শুল্ক} mean? (\textquotedblleft{}A duty; a tax on goods\textquotedblright{}.) Say it once more.
\end{tcolorbox}
\section[\texorpdfstring{\bn{খাজনা} (khājnā)}{খাজনা (khājnā)}]{\bn{খাজনা} (khājnā) — land tax; rent paid on land}

\label{lesson:BN-C418-khajna}



\begin{quote}
Before the new one: say the Bengali for profit (in trade), then the Bengali for a duty (tax).
\end{quote}

\subsection*{You'll want to know: \bn{খাজনা}}
\textbf{\bn{খাজনা}} — \emph{khājnā} — \textquotedblleft{}land tax; rent paid on land\textquotedblright{}.

\begin{quote}
\bn{আগে} \bn{গ্রামের} \bn{লোকেরা} \bn{রাজাকে} \bn{খাজনা} \bn{দিত।}
\end{quote}

\begin{quote}
\emph{āge grāmer lokerā rājāke khājnā ditô.}
\end{quote}

\begin{quote}
In the old days the villagers paid land tax to the king.
\end{quote}

\begin{cousinweb}
\textbf{\bn{খাজনা}} (\emph{khājnā}) is usually traced to Arabic \emph{khazāna}, \textquotedblleft{}a treasury, a store\textquotedblright{}, through Persian: the money that went into the treasury.
\end{cousinweb}

\begin{grammarlens}[title={What you've built}]
One more word for getting around.
\end{grammarlens}

\subsection*{Your turn}
\begin{itemize}
\raggedright
  \item \emph{Say it:} \emph{khājnā}
  \item \emph{Say it:} \emph{khājnā}, once more
  \item \emph{Say it:} \emph{āge grāmer lokerā rājāke khājnā ditô.} — In the old days the villagers paid land tax to the king.
  \item \emph{Recall:} say the Bengali for profit (in trade), then the Bengali for a duty (tax), then say \emph{khājnā} again
\end{itemize}

\begin{tcolorbox}[breakable,colback=teal!4,colframe=teal!35!black,title={Before you move on}]
What does \bn{খাজনা} mean? (\textquotedblleft{}Land tax; rent paid on land\textquotedblright{}.) Say it once more.
\end{tcolorbox}
\section[\texorpdfstring{\bn{মহাজন} (môhājôn)}{মহাজন (môhājôn)}]{\bn{মহাজন} (môhājôn) — a moneylender}

\label{lesson:BN-C418-mohajon}



\begin{quote}
Before the new one: say the Bengali for a duty (tax), then the Bengali for land tax.
\end{quote}

\subsection*{You'll want to know: \bn{মহাজন}}
\textbf{\bn{মহাজন}} — \emph{môhājôn} — \textquotedblleft{}a moneylender\textquotedblright{}.

\begin{quote}
\bn{গরিব} \bn{লোকটা} \bn{মহাজনের} \bn{কাছে} \bn{ধার} \bn{নিল।}
\end{quote}

\begin{quote}
\emph{gôrib lokṭā môhājôner kāchhe dhār nilô.}
\end{quote}

\begin{quote}
The poor man borrowed from the moneylender.
\end{quote}

\begin{cousinweb}
\textbf{\bn{মহাজন}} (\emph{môhājôn}) is Sanskrit \emph{mahā-jana}, \textquotedblleft{}a great or eminent man\textquotedblright{}, also \textquotedblleft{}the head of a trade\textquotedblright{} and, with a query, \textquotedblleft{}a merchant\textquotedblright{} (Monier-Williams); Bengali narrowed it to the one who lends money.
\end{cousinweb}

\begin{grammarlens}[title={What you've built}]
One more word for getting around.
\end{grammarlens}

\subsection*{Your turn}
\begin{itemize}
\raggedright
  \item \emph{Say it:} \emph{môhājôn}
  \item \emph{Say it:} \emph{môhājôn}, once more
  \item \emph{Say it:} \emph{gôrib lokṭā môhājôner kāchhe dhār nilô.} — The poor man borrowed from the moneylender.
  \item \emph{Recall:} say the Bengali for a duty (tax), then the Bengali for land tax, then say \emph{môhājôn} again
\end{itemize}

\begin{tcolorbox}[breakable,colback=teal!4,colframe=teal!35!black,title={Before you move on}]
What does \bn{মহাজন} mean? (\textquotedblleft{}A moneylender\textquotedblright{}.) Say it once more.
\end{tcolorbox}
\section[\texorpdfstring{\bn{বাকি} (bāki)}{বাকি (bāki)}]{\bn{বাকি} (bāki) — the rest; what is left; still owing}

\label{lesson:BN-C418-baki}



\begin{quote}
Before the new one: say the Bengali for land tax, then the Bengali for a moneylender.
\end{quote}

\subsection*{You'll want to know: \bn{বাকি}}
\textbf{\bn{বাকি}} — \emph{bāki} — \textquotedblleft{}the rest; what is left; still owing\textquotedblright{}.

\begin{quote}
\bn{বিলের} \bn{বাকি} \bn{টাকা} \bn{কাল} \bn{দেব।}
\end{quote}

\begin{quote}
\emph{biler bāki ṭākā kāl debô.}
\end{quote}

\begin{quote}
I will pay the rest of the bill tomorrow.
\end{quote}

Also on credit: \bn{বাকিতে} \bn{কেনা}, to buy on credit.

\begin{cousinweb}
\textbf{\bn{বাকি}} (\emph{bāki}) is usually traced to Arabic \emph{bāqī}, \textquotedblleft{}remaining, left over\textquotedblright{}, which reached Bengali through Persian.
\end{cousinweb}

\begin{grammarlens}[title={What you've built}]
One more word for getting around.
\end{grammarlens}

\subsection*{Your turn}
\begin{itemize}
\raggedright
  \item \emph{Say it:} \emph{bāki}
  \item \emph{Say it:} \emph{bāki}, once more
  \item \emph{Say it:} \emph{biler bāki ṭākā kāl debô.} — I will pay the rest of the bill tomorrow.
  \item \emph{Recall:} say the Bengali for land tax, then the Bengali for a moneylender, then say \emph{bāki} again
\end{itemize}

\begin{tcolorbox}[breakable,colback=teal!4,colframe=teal!35!black,title={Before you move on}]
What does \bn{বাকি} mean? (\textquotedblleft{}The rest; what is left; still owing\textquotedblright{}.) Say it once more.
\end{tcolorbox}
